\section{Epistemic and Other Reasoning Modes}\label{sec:epistemic}

The substrate generalizes beyond what is \emph{true} or \emph{probable}. This section covers three further modes that reuse the same CUDA join, circuit, and CDCL machinery---epistemic reasoning over world views, probabilistic aggregates, and well-founded semantics---broadening what a neural network can be integrated with while keeping every mode device-resident.

\subsection{Modal Surface and the Epistemic IR}

Epistemic logic programming asks what is \emph{known} or \emph{possible} across the multiple stable models a program admits---the natural setting for introspection and reasoning under incomplete knowledge. A program may use four modal literals over a predicate---\texttt{know}, \texttt{possible}, \texttt{not know}, \texttt{not possible}. Rather than rewrite them away at parse time, the compiler preserves them in an Epistemic IR (EIR) between the AST and GPU lowering, so the world-view semantics are resolved by a dedicated plan rather than smuggled into relational algebra. The accepted surface spans finite nested modal chains, epistemic integrity constraints (which prune candidate world views on the GPU), rules that couple several epistemic predicates, same-name multi-arity modal predicates, and recursive epistemic execution.

\subsection{World-View Semantics}

A world view is a non-empty set of stable models. \texttt{know p} holds when \texttt{p} appears in every model of the world view; \texttt{possible p} holds when it appears in at least one. Two semantics are selectable: the default, FAEEL (Founded Autoepistemic Equilibrium Logic), rejects circular epistemic support unless it is independently founded, so a self-supporting rule such as \texttt{p :- possible p.} fails closed; a Gelfond-1991-style compatibility mode admits such self-support for legacy specifications. Non-epistemic programs evaluate identically under both.

\subsection{GPU-Native Execution}\label{subsec:epistemic-gpu}

Accepted programs lower from EIR to a GPU plan following a Generate--Propagate--Test schedule whose phases run as CUDA kernels: \emph{generate} enumerates bounded candidate assumptions as device bitsets; \emph{propagate} prunes candidates that contradict known or rejected facts; and \emph{test} validates survivors against the program's world views, checking stable-model tuple membership on-device per arity. The reduced ordinary part of each rule reuses the optimizer, WCOJ planner, and helper-splitting machinery of \Cref{sec:datalog}, and \emph{epistemic splitting} decomposes a program into independent components solved separately and recombined. Candidate validation reuses the on-GPU CDCL solver already built for compilation verification (\Cref{sec:prob})---the same device-resident substrate handles satisfiability, bounded MaxSAT, and assumption-based queries---and accepted world-view evidence feeds the existing GPU-native exact path, so a query can combine epistemic and probabilistic structure without leaving the device. The accepted hot path performs no CPU fallback; a transfer budget tracks that no data-plane host transfers occur within it.

\subsection{Probabilistic Aggregates}

Finite aggregate queries (\texttt{count}, \texttt{sum}, \texttt{min}, \texttt{max}, \texttt{logsumexp}) are supported in both exact provenance/PIR inference and Monte Carlo execution, subject to a typed exact-domain cap that rejects domains too large for tractable enumeration. For small finite \texttt{count} domains, \emph{aggregate lifting} replaces naive world enumeration with exact cardinality dynamic programming over a compact state set---collapsing, for example, a $2^{17}$-outcome enumeration into a few hundred lifted states with identical results---and accepted \texttt{count} aggregates dispatch a GPU-native exact evaluator rather than a CPU finite-world shortcut, keeping the path device-resident.

\subsection{Well-Founded Semantics}

Programs that violate stratification---such as the classic \texttt{p :- not q.\ q :- not p.}---have no two-valued stable model and are rejected by most Datalog engines. xlog instead assigns three-valued truth (true, false, undefined) via Well-Founded Semantics, computing an alternating fixed point: mark the greatest unfounded set false, derive consequences true, and repeat until stable. For probabilistic programs, true atoms receive normal probability and gradient while false and undefined atoms are assigned probability zero, matching the conservative treatment of non-stratified programs. WFS is invoked automatically when a non-monotone SCC is detected during provenance extraction.
